\section{Conclusion}

MetaShift-Bench is a reproducible protocol for auditing reported PM2.5
measurement-method transitions with metadata anchors, distinct physical donor
controls, stable-window synthetic truth, and complete real-event accounting.
Its evidence supports a benchmark-and-audit contribution: it records where
cross-site diagnostics are available, where they are not, and what the
available diagnostics can and cannot say.

The frozen results do not establish a confidence-supported aggregate MetaShift
advantage over standard synthetic control. They also do not establish that a
reported Method Code transition represents a physical instrument change or a
measurement bias. The appropriate use is therefore to prioritize a limited set
of candidate station-specific discontinuities for additional records-based and
technical review, while preserving unsupported and inconclusive anchors in the
public audit trail.

\begin{table}[!ht]
\centering
\caption{Frozen evidence release verification record.}
\label{tab:reproducibility}
\small
\begin{tabular}{lr}
\toprule
Record & Value \\
\midrule
Evidence tag & v0.3.2-evidence-final \\
Frozen evidence commit & \texttt{57d678ecabeb} \\
Synthetic result label & stable\_full\_v2 \\
Case manifest SHA-256 & \texttt{065b1b65c231c529\ldots} \\
Release-gate checks passed & 35/35 \\
Document-consistency checks passed & 12/12 \\
Markdown number checks passed & 57/57 \\
Two-environment core hashes matched & 34 \\
\bottomrule
\end{tabular}
\par\smallskip
\begin{minipage}{0.94\linewidth}
\footnotesize\textit{Note.} The two-environment statement applies only to the listed designated core artifacts in the frozen reproducibility comparison, not to every local file.
\end{minipage}
\end{table}

